\documentclass[
    language=english,
    doctype=proceedings,
    institution=none,
]{../../omnilatex}

\RequirePackage{config/document-settings}
\addbibresource{../../bib/bibliography.bib}

\begin{document}

\maketitle

\section*{Preface}
This volume contains the proceedings of the 8th International Conference on
Sustainable Energy Systems (ICSES 2026), held in Hannover, Germany from
15--17 June 2026.  The conference brought together 180 researchers and
practitioners from 32 countries to present advances in renewable energy
technologies, grid integration, and energy policy~\cite{bibid}.  Of 245
submitted manuscripts, 78 were selected for oral presentation and 94 for poster
display.
This volume publishes extended versions of 42 selected papers, each reviewed by
at least two independent experts.

\tableofcontents

\section{Conference Overview}

\subsection{Technical Programme}
The technical programme was organised into six parallel tracks:
\begin{itemize}
    \item Photovoltaic Systems and Materials
    \item Wind Energy Integration
    \item Energy Storage Technologies
    \item Smart Grid Optimisation
    \item Energy Policy and Economics
    \item Hydrogen and Fuel Cells
\end{itemize}

\subsection{Invited Speakers}
The conference featured keynote lectures from leading researchers in sustainable
energy~\cite{einstein1905}, including presentations on perovskite solar cell
stability, off-shore wind farm control, and sector-coupling
strategies~\cite{goossensLaTeXCompanion1993}.

\section{Selected Papers}

\subsection{Optimal Sizing of Battery Energy Storage Systems for Microgrids}
\textbf{Authors:} H.\ Martinez, S.\ Kim, and R.\ Becker

This paper presents a mixed-integer linear programming formulation for the
optimal sizing of battery energy storage systems (BESS) in islanded
microgrids.  The proposed model minimises total annualised cost while
satisfying reliability constraints defined by the system average interruption
duration index (SAIDI).  Case studies on a rural electrification project in
Sub-Saharan Africa demonstrate cost reductions of 18\% compared to
rule-of-thumb sizing methods.

\subsection{Data-Driven Demand Response in Smart Distribution Networks}
\textbf{Authors:} L.\ Tanaka and P.\ Okonkwo

We propose a data-driven demand response framework that uses recurrent neural
networks to forecast household-level load profiles and aggregate flexibility
resources.  The approach achieves a mean absolute percentage error below 4\%
for day-ahead load forecasting and enables real-time demand response events
with a response latency of less than 30 seconds.

\subsection{Green Hydrogen Production via PEM Electrolysis: A Techno-Economic Assessment}
\textbf{Authors:} F.\ Al-Rashid and M.\ Johansson

This study evaluates the levelised cost of hydrogen (LCOH) for proton exchange
membrane (PEM) electrolysis systems paired with dedicated solar photovoltaic
arrays.  The analysis considers five locations across North Africa and the
Middle East, accounting for variations in solar irradiance, electricity price,
and water availability.

\subsection{Topology Optimisation for Resilient Power Distribution Networks}
\textbf{Authors:} K.\ Novak and T.\ Andersen

We develop a graph-based topology optimisation method for power distribution
networks that enhances resilience against extreme weather events.  The proposed
method formulates the problem as a stochastic mixed-integer second-order cone
program and solves it using a Benders decomposition approach.

\section{Acknowledgements}
The organisers thank the scientific committee, session chairs, and reviewers for
their invaluable contributions.  Financial support from the German Research
Foundation (DFG) and the European Union Horizon Programme is gratefully
acknowledged.

\printbibliography

\end{document}
